% Part: computability
% Chapter: recursive-functions
% Section: bounded-minimization

\documentclass[../../../include/open-logic-section]{subfiles}

\begin{document}

\olfileid[tr]{cmp}{rec}{bmi}
\olsection{Sınırlı En Küçükleme}

\begin{explain}
Bir fonksiyonu, belirli bir $P$ özelliğini ya da bağıntısını sağlayan en küçük
sayı olarak tanımlamak çoğu kez yararlıdır. $P$ karar verilebilirse bu
fonksiyonu, $P$'yi sağlayan ilk sayıyı buluncaya kadar $0,1,2,\dots$
sayılarının tümünü deneyerek hesaplayabiliriz. Bu tür sınırsız bir arama bizi
ilkel özyinelemeli fonksiyonların dışına çıkarır. Fakat yalnızca
\emph{bağımsız olarak verilmiş bir sınırdan küçük} en küçük sayıyla
ilgileniyorsak ilkel özyinelemeli alanda kalırız. Biraz daha genel söylersek,
ilkel özyinelemeli bir $R(x,z)$ bağıntımız olsun. $x$ ile $y$'yi,
$R(x,z)$ koşulunu sağlayan en küçük $z<y$ sayısına götüren fonksiyonu ele
alalım. Bu da $R(x,0),R(x,1),\dots,R(x,y-1)$ koşullarını sınayarak
hesaplanabilir. Peki neden ilkel özyinelemelidir?
\end{explain}

\begin{prop}
$R(\vec x,z)$ ilkel özyinelemeliyse, böyle bir sayı varsa $R(\vec x,z)$
koşulunu sağlayan en küçük $z<y$ sayısını, yoksa $y$'yi döndüren
$m_R(\vec x,y)$ fonksiyonu da ilkel özyinelemelidir. $m_R$ fonksiyonunu
\[
  \bmin{z<y}{R(\vec x,z)}
\]
biçiminde yazacağız.
\end{prop}

\begin{proof}
Hiçbir $z<0$ bulunmadığından $R(\vec x,z)$ koşulunu sağlayan bir $z<0$ da
bulunamaz. Dolayısıyla $m_R(\vec x,0)=0$ olur.

Sınır $y+1$ biçimindeyse üç durum vardır:
\begin{enumerate}
\item $R(\vec x,z)$ koşulunu sağlayan bir $z<y$ vardır. Bu durumda
  $m_R(\vec x,y+1)=m_R(\vec x,y)$.
\item Böyle bir $z<y$ yoktur, fakat $R(\vec x,y)$ geçerlidir. Bu durumda
  $m_R(\vec x,y+1)=y$.
\item $R(\vec x,z)$ koşulunu sağlayan hiçbir $z<y+1$ yoktur. Bu durumda
  $m_R(\vec x,y+1)=y+1$.
\end{enumerate}
Böylece $m_R$ fonksiyonunu ilkel özyinelemeyle şöyle tanımlayabiliriz:
\begin{align*}
  m_R(\vec x,0) & = 0\\
  m_R(\vec x,y+1) & = \begin{cases}
    m_R(\vec x,y) & \text{$m_R(\vec x,y)\neq y$ ise}\\
    y & \text{$m_R(\vec x,y)=y$ ve $R(\vec x,y)$ ise}\\
    y+1 & \text{aksi durumda.}
  \end{cases}
\end{align*}
$R(\vec x,z)$ koşulunu sağlayan bir $z<y$ bulunması ancak ve ancak
$m_R(\vec x,y)\neq y$ olduğunda geçerlidir.
\end{proof}

\begin{prob}
$R(\vec x,z)$ ilkel özyinelemeli olsun. Böyle bir sayı varsa
$R(\vec x,z)$ koşulunu sağlayan en küçük $z<y$ sayısını, yoksa $0$'ı
döndüren $m'_R(\vec x,y)$ fonksiyonunu $\Char{R}$'den ilkel özyinelemeyle
tanımlayın.
\end{prob}

\end{document}

